\chapter[Introdução]{Introdução}

A convergência de competências multidisciplinares é o alicerce da disciplina de Projeto Integrador 1 (FGA0303), que propõe aos estudantes a resolução de problemas complexos por meio da engenharia prática. O presente trabalho detalha a concepção e implementação do sistema micromouse, projetado para navegação autônoma em ambientes labirínticos. O desafio exige a sincronização precisa entre algoritmos de tomada de decisão, processamento de sinais de sensores e uma estrutura mecânica otimizada para agilidade e estabilidade.

Historicamente, a robótica móvel autônoma encontrou seu espaço em soluções industriais de alto desempenho, porém, o acesso a tais tecnologias é frequentemente limitado por custos elevados de implementação e licenciamento. Nesse cenário, o desenvolvimento de protótipos baseados em sistemas embarcados de baixo custo, como o microcontrolador ESP32, surge como uma alternativa viável e educativa. O projeto foca em demonstrar que a eficiência na exploração de labirintos não depende exclusivamente de \textit{hardware} oneroso, mas sim da integração inteligente entre sensores ultrassônicos e controle em malha fechada.

Sob a ótica normativa, embora a legislação brasileira ainda não possua diretrizes específicas para a circulação de robôs autônomos de pequeno porte, a segurança operacional é balizada pela norma regulamentadora NR-12. Esta norma oferece os fundamentos necessários para mitigar riscos técnicos e garantir a integridade do equipamento e dos operadores. Adicionalmente, padrões internacionais como a ANSI/ITSDF B56.5-2019 e a EN ISO 3691-4:2020 servem de guia para o desenvolvimento de sistemas de segurança e navegação estáveis, elevando o rigor técnico do projeto para além do ambiente acadêmico.

A motivação para o desenvolvimento do micromouse fundamenta-se na busca por uma plataforma de telemetria superior às soluções de mercado. Ao contrário de modelos que operam de forma isolada, este projeto introduz uma interface de monitoramento em tempo real (Mission Control), capaz de registrar variáveis cinemáticas e de consumo energético via protocolo MQTT. Esta abordagem não apenas resolve o problema físico do labirinto, mas também melhora a análise de dados do percurso, justificando sua aplicação como um sistema de baixo custo, escalável e tecnicamente robusto para o estudo da robótica autônoma.